\chapter{実験2：洗浄・除去系技術を対象とした異質解決策型補助分析}

\section{本章の概要}
\begin{figure}[tbp]
\centering
\includegraphics[width=0.92\linewidth]{experiment2_pipeline}
\caption{実験2の処理フロー。S1用途クラスタ、C1非半導体系pre候補、各クラスタ上位20件、7クラスタ分、ファミリー単位の重複統合、46件のLLM支援付き著者確認評価を示す。}
\label{fig:experiment2_pipeline}
\end{figure}



本章では、卒論の補助分析として、洗浄・除去系技術から半導体洗浄に接続しうる異質解決策型の専門家確認候補を抽出できるかを確認する。第5章では、U-Netを対象とした近接解決策型の主分析を行った。これに対して本章では、対象物や工程がより異なる場合に、課題文脈を手がかりとして候補を整理できるかを検討する。

本章で扱う実験2は、技術導入可能性や企業行動を示すものではない。花王の洗浄・除去系特許を用いるが、これは企業の将来行動を論じるためではなく、洗浄・除去という課題文脈を持つ特許集合から、半導体洗浄課題との対応を専門家が確認すべき候補として整理できるかを確認するためである。

本章では、まず実験2の目的を述べ、次にS1用途クラスタ、C1候補母集団、クラスタ別候補ランキングの作成手順を説明する。その後、46件のユニークファミリーを対象としたLLM支援付き著者確認評価の結果を示し、最後に軽量全文テキストベースラインとの補助比較を行う。

\section{実験2の目的}

本節では、実験2の目的と位置づけを説明する。実験2は、本卒論の補助分析であり、洗浄・除去系技術から半導体洗浄への異質解決策型候補抽出を検討する。

実験1では、医療画像解析と製造欠陥検出のように、課題と解決手段の両方に比較的近さがある設定を扱った。一方、実験2では、洗浄・除去という課題には一定の関連があるものの、対象物、工程、材料、利用分野が大きく異なる候補を扱う。このような設定では、全文類似だけでは候補抽出の理由が説明しにくく、課題文脈と解決手段文脈を分けて整理する必要がある。

本実験の目的は、半導体洗浄課題との接続が読み取れる専門家確認候補を抽出できるかを確認することである。したがって、本実験は、特定企業の半導体分野への進出や、候補技術の実装可能性を示すものではない。あくまで、既知事例に関連する候補集合を人間が読んだときに、半導体洗浄課題との対応を解釈できるかを確認する補助分析である。

なお、実験2は、実運用における未知候補推薦そのものを再現したものではなく、既知事例を用いた後ろ向き評価として位置づける。実運用では、対象分野の特許集合から課題文脈を抽出し、課題クラスタを作成したうえで、他分野の候補を探索することが想定される。一方、本実験では、花王の半導体洗浄関連既知事例であるS1を用途クラスタに整理し、そのクラスタに対応する課題文脈に対して、C1の非半導体系洗浄・除去技術から候補を抽出した。したがって、本実験の結果は企業進出予測ではなく、異質解決策型候補抽出が専門家確認候補を提示しうるかを確認する補助分析として解釈する。

\section{S1用途クラスタの構成}

本節では、半導体洗浄に関する既知事例であるS1を、用途クラスタに分類する方法を説明する。実験2では、S1を半導体洗浄側の参照集合として用い、その用途に基づいて7つのクラスタを作成した。

本研究で用いた7つの用途クラスタは、以下の通りである。

\begin{itemize}
\item 接着剤洗浄クラスタ
\item 回路基板洗浄クラスタ
\item CMP後洗浄クラスタ
\item フラックス・電子部品洗浄クラスタ
\item フォトレジスト・樹脂除去クラスタ
\item 半導体基板洗浄クラスタ
\item シリコンウェーハ洗浄クラスタ
\end{itemize}

各クラスタは、半導体洗浄や電子部品洗浄に関連する用途の違いを表す。例えば、シリコンウェーハ洗浄、半導体基板洗浄、CMPまたは研磨後洗浄、フラックス残渣や電子部品洗浄などに対応する。

各クラスタについて、クラスタ内のS1ファミリーのうち最も早い日付をクラスタ初出日として設定した。このクラスタ初出日は、後続のC1候補抽出において、既知事例より前に存在する候補を抽出するための時系列条件として用いる。

このクラスタ化は、実験2の候補抽出を用途別に比較するための設計である。クラスタごとに課題や対象物が異なるため、候補の解釈可能性にも差が生じる可能性がある。

\section{C1候補母集団の作成}

本節では、実験2におけるC1候補母集団の作成方法を説明する。C1は、花王の洗浄・除去系技術に関する特許集合のうち、半導体直接用途ではない先行候補を抽出するために用いる。

まず、S1と同一ファミリーに該当するものは除外した。これは、既知の半導体洗浄事例そのものが候補集合に混入することを避けるためである。次に、半導体直接語を含むものを除外した。実験2の目的は、半導体用途として既に明示されている特許を再発見することではなく、非半導体系の洗浄・除去技術から、半導体洗浄課題に接続しうる専門家確認候補を抽出することであるためである。

ただし、\texttt{substrate} のような語は、単独では半導体用途を直接示すとは限らない。そのため、\texttt{substrate} 単独の出現は除外条件とはせず、確認用フラグとして扱った。これにより、基板という語を含むが半導体直接用途とは限らない候補を、人手確認の対象として残すことができる。

また、各S1クラスタの初出日を基準として、それ以前に存在するC1候補を抽出した。これにより、各クラスタの既知事例より前に存在する非半導体系の候補を、クラスタ別にランキングする準備を行った。

このC1候補母集団は、専門家確認候補を作るための前処理済み集合であり、候補の妥当性をこの段階で判断するものではない。候補が半導体洗浄課題に接続しうるかどうかは、後続のランキングと人手評価により確認する。

\section{クラスタ別候補ランキング}

本節では、S1用途クラスタごとにC1候補をランキングする方法を説明する。実験2では、7クラスタそれぞれについて、課題文脈に基づく候補ランキングを作成した。

まず、各クラスタに対して、クラスタ初出日以前に存在するS0側の比較用集合を作成した。次に、用途語フィルタを用いて、各クラスタの課題に対応しやすいS0\_pre\_filteredを作成した。これにより、クラスタごとの課題文脈をより明確にし、その文脈をもとにC1候補をランキングできるようにした。

各S1用途クラスタについて、クラスタに対応するS0\_pre\_filteredの\texttt{problem\_context}を収集し、これを対象分野側の課題文脈として扱った。次に、これらの\texttt{problem\_context}をBERT系モデルによりベクトル化し、クラスタごとの課題文脈を表すベクトル集合を作成した。さらに、C1候補の\texttt{problem\_context}も同様にベクトル化し、両者の類似度に基づいてクラスタ別ランキングを作成した。これにより、半導体用途そのものではないが、半導体洗浄課題と類似した課題を扱う洗浄・除去系特許を専門家確認候補として抽出した。

なお、このランキングは候補の技術導入可能性を判定するものではない。上位候補については、\texttt{solution\_context}、\texttt{technical\_means}、\texttt{target\_object}を確認し、半導体洗浄課題との対応が解釈可能かをLLM支援付き著者確認評価により確認した。

その後、各クラスタについてC1候補をランキングし、上位20件の候補を抽出した。7クラスタそれぞれで上位20件を抽出するため、各クラスタ上位20件、7クラスタ分の候補が得られる。ただし、同一ファミリーが複数クラスタに出現する場合があるため、最終的にはファミリー単位で重複統合を行った。

重複統合の結果、46件のユニークファミリーをLLM支援付き著者確認評価の対象とした。この46ファミリーは、提案手法により抽出された実験2の評価対象である。ただし、この時点では候補の妥当性を確定するものではなく、専門家確認候補として人間が読むための候補集合である。

\section{LLM支援付き著者確認評価}
\begin{figure}[tbp]
\centering
\includegraphics[width=0.92\linewidth]{experiment2_label_distribution}
\caption{実験2における評価ラベルの分布。46件のユニークファミリーを対象にしたLLM支援付き著者確認評価の結果を示す。}
\label{fig:experiment2_label_distribution}
\end{figure}

\begin{figure}[tbp]
\centering
\includegraphics[width=0.92\linewidth]{experiment2_cluster_distribution}
\caption{実験2におけるクラスタ別評価ラベル分布。クラスタ依存性を確認するための補助図である。}
\label{fig:experiment2_cluster_distribution}
\end{figure}

\begin{table}[tbp]
\centering
\small
\caption{実験2におけるLLM支援付き著者確認評価の全体集計}
\label{tab:experiment2_label_summary_table}
\begin{adjustbox}{max width=\textwidth}
\begin{tabular}{lll}
\toprule
ラベル & 意味 & 件数 \\
\midrule
◎ & 強い候補 & 11 \\
○ & 妥当候補 & 5 \\
△ & 弱い候補 & 12 \\
× & 不適切候補 & 18 \\
◎+○ & 強い候補または妥当候補 & 16 \\
◎+○+△ & 広義の確認候補 & 28 \\
\bottomrule
\end{tabular}
\end{adjustbox}
\end{table}


\begin{table}[tbp]
\centering
\small
\caption{実験2におけるクラスタ別評価集計}
\label{tab:experiment2_cluster_summary_table}
\begin{adjustbox}{max width=\textwidth}
\begin{tabular}{llllllll}
\toprule
クラスタ & N & ◎ & ○ & △ & × & ◎+○ & ◎+○+△ \\
\midrule
接着剤洗浄 & 7 & 2 & 0 & 0 & 5 & 2 & 2 \\
回路基板洗浄 & 3 & 1 & 1 & 1 & 0 & 2 & 3 \\
CMP後洗浄 & 12 & 0 & 0 & 5 & 7 & 0 & 5 \\
フラックス・電子部品洗浄 & 5 & 1 & 1 & 2 & 1 & 2 & 4 \\
フォトレジスト・樹脂除去 & 6 & 2 & 0 & 2 & 2 & 2 & 4 \\
半導体基板洗浄 & 6 & 5 & 1 & 0 & 0 & 6 & 6 \\
シリコンウェーハ洗浄 & 7 & 0 & 2 & 2 & 3 & 2 & 4 \\
\bottomrule
\end{tabular}
\end{adjustbox}
\end{table}




本節では、実験2で抽出した46件のユニークファミリーに対するLLM支援付き著者確認評価の結果を示す。評価では、候補が半導体洗浄課題との対応を専門家確認候補として解釈できるかを確認した。

評価ラベルは、\texttt{◎}、\texttt{○}、\texttt{△}、\texttt{×} の4段階とした。\texttt{◎} は半導体洗浄、基板洗浄、残渣除去などに直接接続しうる強い候補を表す。\texttt{○} は精密洗浄、電子部品、材料表面処理など、周辺領域として接続しうる候補を表す。\texttt{△} は課題や効果は近いが、対象物、工程、解決手段の対応が弱い候補を表す。\texttt{×} は、異質解決策候補としての対応が読み取りにくい候補を表す。

\figref{fig:experiment2_label_distribution}および\tabref{tab:experiment2_label_summary_table}に、最終評価の全体分布を示す。46件のうち、\texttt{◎} は11件、\texttt{○} は5件、\texttt{△} は12件、\texttt{×} は18件であった。\texttt{◎+○} は16件であり、強い候補または妥当な専門家確認候補として扱える候補が一定数得られた。また、\texttt{◎+○+△} は28件であり、広義には確認対象として残しうる候補も含まれている。

一方で、\texttt{×} も18件含まれており、全体としてノイズは多い。これは、異質解決策型の候補抽出では、洗浄・除去という課題が近くても、対象物や汚れの種類、工程条件が異なるため、半導体洗浄課題との対応が読み取りにくい候補が混入しやすいためである。

\figref{fig:experiment2_cluster_distribution}および\tabref{tab:experiment2_cluster_summary_table}に、クラスタ別の評価結果を示す。半導体基板洗浄クラスタでは6件中5件が\texttt{◎}、1件が\texttt{○}であり、強い候補が多く得られた。これは、基板洗浄という対象や課題が半導体洗浄と比較的近いため、候補の対応関係を読み取りやすかったためと考えられる。

一方、CMP後洗浄クラスタでは12件中\texttt{◎}と\texttt{○}が0件であり、\texttt{△}が5件、\texttt{×}が7件であった。このクラスタでは、CMP後洗浄という工程固有性が強く、非半導体系の洗浄・除去技術から直接対応を読み取ることが難しい候補が多かった。接着剤洗浄クラスタでも7件中\texttt{×}が5件であり、候補の対象物や汚染物の違いが評価に影響した。

これらの結果から、実験2では一部クラスタで専門家確認候補として解釈可能な候補が得られた一方、クラスタによって候補の解釈可能性に大きな差があることが確認された。したがって、異質解決策型の候補抽出では、クラスタ設計や用途語フィルタが結果に強く影響する。

\section{軽量全文テキストベースラインとの補助比較}
\begin{figure}[tbp]
\centering
\includegraphics[width=0.92\linewidth]{experiment2_fulltext_overlap}
\caption{提案手法と軽量全文テキストベースラインの候補集合の重複・差分。精度優劣ではなく候補集合の傾向差を確認する。}
\label{fig:experiment2_fulltext_overlap}
\end{figure}

\begin{table}[tbp]
\centering
\small
\caption{実験2における提案手法と軽量全文テキストベースラインの比較}
\label{tab:experiment2_fulltext_baseline_table}
\begin{adjustbox}{max width=\textwidth}
\begin{tabular}{lll}
\toprule
指標 & 内容 & 値 \\
\midrule
提案手法件数 & 提案手法のファミリー数 & 46 \\
ベースライン件数 & 軽量全文テキストベースラインのファミリー数 & 56 \\
重複件数 & 重複ファミリー数 & 11 \\
提案手法のみ & 提案手法にのみ含まれるファミリー数 & 35 \\
ベースラインのみ & ベースラインにのみ含まれるファミリー数 & 45 \\
\bottomrule
\end{tabular}
\end{adjustbox}
\end{table}




本節では、提案手法による候補集合と、軽量全文テキストベースラインによる候補集合を比較する。ここで用いるベースラインは、厳密なBERT型全文埋め込み実験ではなく、Title、Abstract、Claimsを用いた軽量全文テキストベースラインである。

実験2では、提案手法による課題文脈ベースのランキングと、全文テキストに基づく候補集合がどの程度重なるかを確認した。\tabref{tab:experiment2_fulltext_baseline_table}に、ファミリー単位の重複比較を示す。提案手法の候補数は46ファミリー、軽量全文テキストベースラインの候補数は56ファミリーであった。このうち、重複したファミリーは11件であり、提案手法のみに含まれる候補は35件、軽量全文テキストベースラインのみに含まれる候補は45件であった。

\figref{fig:experiment2_fulltext_overlap}に、提案手法と軽量全文テキストベースラインの重複関係を示す。重複が11件にとどまったことから、課題文脈に基づくランキングと、全文テキスト類似に基づくランキングは、異なる候補集合を提示することが分かる。これは、全文テキストでは対象物や用途語、請求項中の表現が強く影響する一方で、課題文脈ベースのランキングでは「何の課題に対応しているか」に焦点を当てるためである。

重複した11件の内訳は、\texttt{◎}が4件、\texttt{○}が3件、\texttt{△}が2件、\texttt{×}が2件であった。この内訳から、両方の方法に共通して現れる候補には一定数の強い候補や妥当候補が含まれる一方、不適切候補も含まれることが分かる。

ただし、軽量全文テキストベースラインのみに含まれる候補は未評価である。そのため、本比較から、提案手法がベースラインより精度面で優れているとは結論づけない。本比較は、候補集合の傾向差を確認するための補助分析である。提案手法の意義は、全文テキスト類似とは異なる観点から、専門家確認候補を抽出し、その抽出理由を課題文脈に基づいて説明しやすくする点にある。

\section{実験2のまとめ}

本節では、実験2の結果を補助分析としてまとめる。実験2では、洗浄・除去系技術から半導体洗浄に接続しうる異質解決策型の専門家確認候補を抽出できるかを確認した。

まず、S1を7つの用途クラスタに分類し、各クラスタについてC1非半導体系の先行候補をランキングした。各クラスタ上位20件、7クラスタ分の候補をファミリー単位で重複統合し、46件のユニークファミリーをLLM支援付き著者確認評価の対象とした。

最終評価の結果、\texttt{◎}が11件、\texttt{○}が5件、\texttt{△}が12件、\texttt{×}が18件であった。\texttt{◎+○}は16件、\texttt{◎+○+△}は28件であり、全体としてノイズは多いものの、一部には専門家確認候補として解釈可能な候補が得られた。特に、半導体基板洗浄クラスタでは強い候補が多く、CMP後洗浄クラスタでは強い候補が少ないなど、クラスタ依存性が確認された。このことから、異質解決策型の候補抽出では、対象分野側の課題クラスタ設計が候補品質に大きく影響することが示唆された。

また、軽量全文テキストベースラインとの補助比較では、提案手法とベースラインの重複は11ファミリーにとどまった。これは、課題文脈ベースのランキングと全文テキストベースのランキングが、異なる候補集合を提示することを示している。ただし、ベースラインのみの候補は未評価であるため、本比較を手法間の精度優劣として扱うことはできない。

以上より、実験2は、異質解決策型においても、課題文脈を手がかりに専門家確認候補を整理できる可能性を補助的に示した。一方で、候補の解釈可能性はクラスタに大きく依存し、ノイズも多い。したがって、実験2は本卒論における主結果ではなく、提案手法の適用範囲を探索的に確認する補助分析として位置づける。
